\section{子任务 10：第 9 章 Intractability}

\subsection{核心知识点}

\begin{itemize}[leftmargin=2em]
  \item hierarchy theorems：更多时间或空间确实带来更强计算能力。
  \item 指数时间、指数空间中的完备性现象。
  \item relativization：给机器加 oracle 后，复杂度类之间的关系可能改变。
  \item 相对化结果揭示某些证明方法的局限。
  \item circuit complexity：用非一致模型研究计算难度。
\end{itemize}

\subsection{典型问题与证明套路}

\begin{enumerate}[leftmargin=2em]
  \item \textbf{证明层级严格}：做受资源约束的对角化，让新机器在给定资源内专门和低层机器作对。
  \item \textbf{证明某方法不够强}：构造 oracle 世界，使得同一方法在某些世界里支持结论、另一些世界里反对结论。
  \item \textbf{讨论电路复杂度}：把算法问题转成布尔函数族，看需要多大电路、多少层深、何种非一致资源。
\end{enumerate}

\subsection{证明背后的哲学}

这一章的哲学非常成熟：\textbf{不仅要证明结论，还要反思证明工具本身。}
层级定理说明资源差异确实重要；相对化则提醒我们，某些自然方法即使强大，也可能永远不足以解决像 $P$ vs.\ $NP$ 这样的核心难题。理论开始从“研究问题”上升到“研究证明问题的方法论”。

\subsection{本章精华}

\begin{itemize}[leftmargin=2em]
  \item Intractability 不只是说“这题难”，而是研究“为什么我们现有工具难以证明它不难”。
  \item 复杂性理论在这一章显出元理论气质：它开始分析自身的证明边界。
\end{itemize}

\newpage
